\begin{table}[pos=htbp]
\centering\small
\caption{Fields of the encoded traveler state. Mode availability is represented separately from these attributes.}
\label{tab:input_dictionary}
\renewcommand{\arraystretch}{1.12}
\begin{tabularx}{\linewidth}{@{}p{0.25\linewidth}X@{}}
\toprule
Input group & Fields \\
\midrule
Categorical traveler fields & age group, income group, occupation, habitual mode, schedule flexibility, mobility limitation \\
Categorical trip/context fields & trip purpose, time constraint, weather condition \\
Numerical traveler fields & household size, presence of children, car ownership, driving licence, bicycle ownership, transit pass \\
Numerical trip fields & distance (km), desired departure time, desired arrival time \\
Numerical context fields & weather intensity, road congestion, PT delay, PT disruption, road disruption, fare multiplier, parking cost multiplier, congestion charge \\
Alternative fields & travel time, monetary cost, access time, number of transfers, reliability delay, weather exposure \\
Additional PT-supply fields & PT feasibility, egress time, waiting time, in-vehicle time, transfer time, coverage ratio \\
\bottomrule
\end{tabularx}
\par\smallskip\begin{minipage}{\linewidth}\footnotesize
The nine categorical and seventeen global numerical fields form the shared state; each mode has twelve numerical attributes. Times are in minutes and distance in kilometres before encoding. Categorical vocabularies and numerical normalization are fitted on training data. For SA-Student, existing statistics are inherited and only the six added supply fields receive newly fitted normalization. The supplementary data files give the variable names used in the implementation.
\end{minipage}
\end{table}
